\documentclass[aspectratio=169]{beamer}

\usepackage[T2A]{fontenc}
\usepackage[utf8]{inputenc}
\usepackage[russian,english]{babel}
\usepackage{cmap}
\usepackage{tikz}
\usepackage{array}
\usepackage{booktabs}
\usepackage{graphicx}
\usepackage{csquotes}
\usepackage{amssymb,amsfonts,amsmath,mathtools}

\usetikzlibrary{arrows.meta,positioning}

\usecolortheme{beaver}
\usefonttheme[onlymath]{serif}
\setbeamertemplate{navigation symbols}{}
\setbeamertemplate{footline}[page number]
\setbeamertemplate{blocks}[rounded=true,shadow=true]
\setbeamersize{text margin left=0.55cm,text margin right=0.55cm}
\setbeamerfont{author}{size=\normalsize}
\setbeamerfont{institute}{size=\normalsize}
\setbeamerfont{date}{size=\normalsize}

\newcommand{\model}{Attentive model}
\newcommand{\contextmodel}{Context model}
\newcommand{\plsa}{PLSA}
\newcommand{\bertopic}{BERTopic}
\newcommand{\Perp}{\operatorname{Perplexity}}
\newcommand{\normop}{\operatorname{norm}}

\title{Параметризация тематических моделей локального контекста}
\author{Matvey Karpov}
\institute{
    Intelligent Systems Department, MIPT
}
\date{\today}

\begin{document}
\selectlanguage{russian}

\begin{frame}
    \thispagestyle{empty}
    \maketitle
\end{frame}

\begin{frame}{Мотивация}
    Классические вероятностные тематические модели, включая \plsa{} и LDA,
    обычно используют представление документа как мешка слов. В такой
    постановке правдоподобие зависит от агрегированных счетчиков термов, но не
    от порядка токенов.

    \begin{columns}[T]
        \begin{column}{0.49\textwidth}
            \begin{block}{Ограничение}
                Документы с одинаковыми счетчиками \(n_{dw}\), но различным
                порядком слов, неразличимы для модели мешка слов.
            \end{block}
        \end{column}
        \begin{column}{0.49\textwidth}
            \begin{block}{Исследуемая гипотеза}
                Локальный контекст \(C_i\) позиции \(i\) может быть включен в
                вероятностную модель и использован при оценивании
                тематического представления токена.
            \end{block}
        \end{column}
    \end{columns}
\end{frame}

\begin{frame}{Цель исследования}
    \small
    \begin{block}{Проблема}
        Модель мешка слов использует счетчики термов, но не учитывает порядок
        токенов и локальные совместные встречаемости слов.
    \end{block}

    \begin{block}{Цель}
        Построить вероятностную тематическую модель, где тема позиции \(i\)
        определяется обучаемо взвешенным локальным контекстом:
        \[
            \theta_{ti}=p(t\mid C_i)
            =\sum_{c\in C_i}\alpha_{ci}\phi_{x_c t},
            \qquad \alpha_i\in\Delta^{|C_i|-1}.
        \]
    \end{block}

    \begin{block}{Нужно}
        \begin{itemize}
            \item вывести обновления для \(\Phi\) и \(\alpha\) при симплексных ограничениях;
            \item сравнить обучаемый контекст с фиксированной \contextmodel{},
            \plsa{} и \bertopic{} по тестовой перплексии и NPMI.
        \end{itemize}
    \end{block}
\end{frame}

\begin{frame}{Постановка}
    После предобработки коллекция задается последовательностью токенов
    \[
        x_1,\ldots,x_N,\qquad x_i\in W,
    \]
    а границы документов задаются индексами
    \(0=b_0<b_1<\cdots<b_{|D|}=N\).

    \begin{block}{Требуется}
        Оценить скрытое тематическое представление, в котором каждое слово и
        каждый локальный контекст сопоставляются распределениям на множестве
        тем \(T\), а строки \(\Phi\) лежат в единичных симплексах.
    \end{block}

    \[
        \Phi=(\phi_{wt})_{w\in W,t\in T},\qquad
        \sum_{t\in T}\phi_{wt}=1,\qquad
        p(t)=\sum_{w\in W}p(w)\phi_{wt}.
    \]

    При \(p(t)>0\) распределение слов в теме восстанавливается по формуле
    Байеса:
    \[
        p(w\mid t)=\frac{p(w)\phi_{wt}}{p(t)}.
    \]
\end{frame}

\begin{frame}{Локальный контекст}
    Для позиции \(i\) локальный контекст \(C_i\) состоит только из индексов
    того же документа. Это исключает перенос контекстного окна через границы
    документов. Каждому элементу контекста сопоставляется коэффициент
    \[
        \alpha_{ci}\ge 0,\qquad \sum_{c\in C_i}\alpha_{ci}=1.
    \]

    \begin{columns}[T]
        \begin{column}{0.53\textwidth}
            \[
                \theta_{ti}=p(t\mid C_i)
                =\sum_{c\in C_i}\alpha_{ci}\phi_{x_c t}.
            \]
            \[
                p(x_i\mid C_i)
                =p(x_i)\sum_{t\in T}
                \frac{\phi_{x_i t}\theta_{ti}}{p(t)}.
            \]
            \[
                p_{ti}=p(t\mid x_i,C_i)
                =\normop_{t\in T}
                \left(\frac{\phi_{x_i t}\theta_{ti}}{p(t)}\right).
            \]
        \end{column}
        \begin{column}{0.45\textwidth}
            \centering
            \begin{tikzpicture}[
                token/.style={draw,rounded corners,minimum width=0.8cm,minimum height=0.48cm},
                arr/.style={-Latex,thick}
            ]
                \node[token] (a) {\(x_{i-2}\)};
                \node[token,right=0.15cm of a] (b) {\(x_{i-1}\)};
                \node[token,right=0.15cm of b,fill=red!10] (c) {\(x_i\)};
                \node[token,right=0.15cm of c] (d) {\(x_{i+1}\)};
                \node[token,right=0.15cm of d] (e) {\(x_{i+2}\)};
                \node[draw,rounded corners,below=1.0cm of c,minimum width=2.8cm,minimum height=0.7cm] (theta) {\(\theta_i=p(t\mid C_i)\)};
                \node[draw,rounded corners,below=0.8cm of theta,minimum width=2.8cm,minimum height=0.7cm] (post) {\(p_i=p(t\mid x_i,C_i)\)};
                \draw[arr] (a) -- (theta);
                \draw[arr] (b) -- (theta);
                \draw[arr] (d) -- (theta);
                \draw[arr] (e) -- (theta);
                \draw[arr] (c) -- (post);
                \draw[arr] (theta) -- (post);
            \end{tikzpicture}
        \end{column}
    \end{columns}
\end{frame}

\begin{frame}{Обучаемые коэффициенты контекста}
    \small
    Коэффициенты \(\alpha_{ci}\) задают вклад слова \(x_c\) из окна \(C_i\) в
    тематическое представление центральной позиции:
    \[
        \theta_{ti}=\sum_{c\in C_i}\alpha_{ci}\phi_{x_c t},
        \qquad
        \alpha_{ci}\ge 0,
        \qquad
        \sum_{c\in C_i}\alpha_{ci}=1.
    \]

    При фиксированных \(\Phi\) и \(p(t)\) локальный блок обучается
    мультипликативным шагом
    \[
        \alpha_{ci}^{(k+1)}
        =
        \normop_{c\in C_i}\left(
        \alpha_{ci}^{(k)}
        \sum_{t\in T}p_{ti}^{(k)}
        \frac{\phi_{x_c t}}{\theta_{ti}^{(k)}}
        \right).
    \]

    То есть вес соседней позиции растет, если ее тематический профиль
    \(\phi_{x_c t}\) хорошо согласуется с апостериорными темами центрального
    токена.
\end{frame}

\begin{frame}{Итерационная процедура}
    \small
    Обучение выполняется как блочная итерационная процедура при симплексных
    ограничениях на строки \(\Phi\) и веса \(a\).

    \[
        \theta_{ti}=\sum_{c\in C_i}\alpha^{(t)}_{ci}\phi_{x_c t},
        \qquad
        p_{ti}=
        \normop_{t\in T}\left(\frac{\phi_{x_i t}\theta_{ti}}{p(t)}\right).
    \]

    \[
        n_{wt}=\sum_i p_{ti}\mathbf{1}[x_i=w],
        \qquad
        N_{wt}=\sum_i
        \left(\sum_{c\in C_i}\alpha_{ci}\mathbf{1}[x_c=w]\right)
        \frac{p_{ti}}{\theta_{ti}}.
    \]

\end{frame}

\begin{frame}{Итерационная процедура: нормировки}
    \small
    \begin{block}{Основные нормировки}
        \[
        \phi_{wt}\leftarrow
        \normop_{t\in T}
        \left(n_{wt}+\phi_{wt}N_{wt}\right),
        \qquad
        p(t)\leftarrow \frac{1}{N}\sum_i p_{ti}.
        \]
        \[
        a_{st}\leftarrow
        \normop_{s}
        \left(
        a_{st}\sum_{i:\,s\in S_i}
        p_{ti}\frac{\phi_{x_{i+s}t}}{\theta_{ti}}
        \right).
        \]
    \end{block}

    
\end{frame}

\begin{frame}{Стационарные условия}
    Для локального вклада
    \[
        \ell_i(\alpha_i)=\log p(x_i\mid C_i)
    \]
    при фиксированных \(\Phi\) и \(p(t)\) производная по коэффициенту
    контекста имеет вид
    \[
        \frac{\partial \ell_i}{\partial \alpha_{ci}}
        =
        \sum_{t\in T}p_{ti}\frac{\phi_{x_c t}}{\theta_{ti}},
        \qquad c\in C_i.
    \]

    \begin{block}{Стационарное уравнение}
        \[
        \alpha_{ci}=
        \normop_{c\in C_i}\left(
        \alpha_{ci}\sum_{t\in T}
        p_{ti}\frac{\phi_{x_c t}}{\theta_{ti}}
        \right).
        \]
    \end{block}

\end{frame}

\begin{frame}{Стационарные условия: сходимость}
    \small
    \begin{block}{Теорема о сходимости локального блока}
        Зафиксируем \(\Phi\), \(p(t)\) и позицию \(i\). Пусть
        \(\alpha_{ci}^{(0)}>0\) для всех \(c\in C_i\) и
        \(p(x_i\mid C_i)>0\) на всех итерациях. Тогда локальные
        мультипликативные обновления не убывают по
        \(\ell_i(\alpha_i)\):
        \[
            \ell_i(\alpha_i^{(k+1)})
            \ge
            \ell_i(\alpha_i^{(k)}).
        \]
        Последовательность \(\ell_i(\alpha_i^{(k)})\) сходится, а
        каждая предельная точка \(\alpha_i^\ast\) удовлетворяет
        KKT-условиям задачи
        \[
            \max_{\alpha_i\in\Delta_i}\ell_i(\alpha_i).
        \]
    \end{block}
\end{frame}

\begin{frame}{Экспериментальный протокол}
    \begin{columns}[T]
        \begin{column}{0.48\textwidth}
            \begin{block}{Корпуса}
                \begin{tabular}{@{}lrr@{}}
                    \toprule
                    Корпус & Train & Test\\
                    \midrule
                    20 Newsgroups & 11314 & 7532\\
                    Yahoo! Answers & 1399999 & 59999\\
                    AG News & 96000 & 7600\\
                    \bottomrule
                \end{tabular}
            \end{block}
        \end{column}
        \begin{column}{0.48\textwidth}
            \begin{block}{Сравнение}
                \begin{itemize}
                    \item \model
                    \item \contextmodel
                    \item \plsa{} в реализации BigARTM
                    \item \bertopic{} с all-MiniLM-L6-v2
                \end{itemize}
            \end{block}
        \end{column}
    \end{columns}

    \vspace{0.25cm}
    Для внешнего сравнения используется не более \(10000\) обучающих и
    \(3000\) тестовых документов на корпус. Основная метрика -- тестовая
    перплексия; дополнительно оцениваются NPMI, время обучения и число
    итераций до остановки.
\end{frame}

\begin{frame}{Сравнение моделей}
    \centering
    \scriptsize
    \begin{tabular}{@{}llrrr@{}}
        \toprule
        Модель & Корпус & Test PPL & NPMI & Время, с\\
        \midrule
        \model & 20News & 383.53 & 0.174 & 29.65\\
        \contextmodel & 20News & 638.54 & 0.177 & 2.58\\
        \plsa & 20News & 4570.90 & 0.252 & 111.00\\
        \bertopic & 20News & 2796.06 & 0.242 & 78.89\\
        \midrule
        \model & Yahoo & 245.93 & 0.097 & 19.18\\
        \contextmodel & Yahoo & 403.82 & 0.100 & 1.58\\
        \plsa & Yahoo & 2851.39 & 0.157 & 67.56\\
        \bertopic & Yahoo & 2600.34 & 0.146 & 70.76\\
        \midrule
        \model & AG News & 294.70 & -0.009 & 8.72\\
        \contextmodel & AG News & 474.93 & -0.015 & 0.94\\
        \plsa & AG News & 16853.97 & 0.129 & 51.52\\
        \bertopic & AG News & 3313.65 & 0.116 & 39.37\\
        \bottomrule
    \end{tabular}

    \vspace{0.18cm}
    \begin{block}{Интерпретация}
        В данном эксперименте \model{} имеет наименьшую тестовую перплексию на
        всех трех корпусах. При этом по NPMI верхних слов темы \plsa{} и
        \bertopic{} остаются сильнее.
    \end{block}
\end{frame}

\begin{frame}{Абляции}
    \begin{columns}[T]
        \begin{column}{0.5\textwidth}
            \begin{block}{Обучение весов контекста}
                Отключение обучения коэффициентов ухудшает перплексию:
                \begin{center}
                \small
                \begin{tabular}{@{}lr@{}}
                    20News & \(383.53\to 518.09\)\\
                    Yahoo & \(245.93\to 597.32\)\\
                    AG News & \(294.70\to 539.13\)
                \end{tabular}
                \end{center}
            \end{block}
            
        \end{column}
        \begin{column}{0.47\textwidth}
            \begin{block}{Радиус окна}
                При фиксированном числе тем широкий диапазон радиусов дает
                близкую перплексию, но большие окна существенно увеличивают
                время обучения.
            \end{block}
            
        \end{column}
    \end{columns}
\end{frame}

\begin{frame}{Выводы}
    \begin{enumerate}
        \item Обучаемые коэффициенты контекста позволяют адаптировать вклад
        соседних позиций; в проведенных экспериментах это уменьшает тестовую
        перплексию относительно фиксированной \contextmodel{}.
        \item Тем не менее, NPMI ниже, чем у \plsa{} и \bertopic{}.
    \end{enumerate}
\end{frame}

\end{document}
